\section{Lista 4}

\setcounter{zadanie}{0}

\begin{lemma}{q.e. z back and forth (3.1.6 David Marker)}{}
  Niech $T$ będzie teorią. Niech $M,N\models T$ będą modelami o wspólnej podstrukturze $A$. Wtedy jeśli dla wszystkich otwartych formuł $\phi(\overline{v}, w)$ zachodzi
  $$(\forall\;\overline{a}\in A)[(\exists\;b\in M)\;M\models \phi(\overline{a}, b)]\implies (\exists\;c\in N)\;N\models\phi(\overline{a}, c)$$
  (czyli chcę używać lematu o rozszerzaniu dla wszystkich $a\in M$ tj. przepchać typ $tp^M(a/A)$ do $tp^N(b/A)$) to teoria $T$ jest q.a..
\end{lemma}

\begin{lemma}{}{}
  T jest q.e. $\iff$ $(\forall\;M,N\models T,\;M,N-\aleph_0-\text{nasycone})\;I(M,N)$ (to są wszystkie częściowe izomorfizmy) ma własność back and forth.
\end{lemma}

to jest chyba na pewno dokładnie to samo co to wyżej

\begin{zadanie}
  Podaj aksjomaty teorii $Th(\omega^\omega, E_n)_{n\in\omega}$, gdzie $E_n(f,g)\iff f(n)=g(n)$, teoria niezależnych relacji równoważności o nieskończenie wielu klasach. Pokaż, że jest ona q.e.
\end{zadanie}

\begin{proof}
  Aksjomaty:
  \begin{itemize}
    \item $E_n$ to relacja równoważności, czyli mamy zwrotność, przechodniość i symetryczność
    \item jest nieskończenie wiele klas relacji $E_n$: 
      $$(\exists\;x_1)\;...\;(\exists\;x_m)\;\bigwedge_{i\neq j} \neg E_n(x_i, x_j)$$
    \item niezależność, czyli dla różnych $n_1$, ..., $n_m$ mamy
      $$(\forall\;x_1)\;...\;(\forall\;x_m)(\exists\;y)\;\bigwedge E_{n_i}(x_i, y)$$
      (w przekrojach zawsze coś jest).
  \end{itemize}

  Q.E.

  Chcę zacząć od częściowego izomorfizmu $f:\{m_1,...,m_k\}\to \{n_1,...,n_k\}$ dla $M,N\models T$. Biorę $a\in M$ i chcę znaleźć $b\in N$ takie, że $tp^M(a/A)=tp^N(b/A)$.

  Jak wygląda $tp^M(a/A)$? Dla każdego $k\in \N$ mam zbiór informacji $\pm E_k(m_i, a)$ dla wszystkich $m_i$. 
  \begin{itemize}
    \item Dla ustalonego $k$ jestem w stanie znaleźć takie $b_k\in N$ które by to spełniało, bo mamy nieskończenie wiele klas relacji $E_n$. 
    \item Teraz mając te wszystkie $b_k$ możemy wziąć element z przekroju klas $b_k$ dzięki aksjomatowi niezależności.
  \end{itemize}
\end{proof}

\subsection{Powtórka: teorie stabilne}

Powiemy, że teoria $T$ ($\mathcal{M}\models T$ jest modelem monstrum w którym pracujemy) jest 
\begin{description}
  \item[$\kappa$-stabilna] jeśli $\kappa\geq \aleph_0$ oraz $(\forall\;A\subseteq M, |A|\leq \kappa)|S_1(A)|\leq \kappa$
  \item[stabilna] gdy jest $\kappa$-stabilna dla pewnego $\kappa\geq\aleph_0$
  \item[superstabilna] gdy istnieje $\mu\geq\aleph_0$ że dla każdego $\kappa\geq\mu$ jest $\kappa$-stabilna
\end{description}

Rozważamy poniższe teorie
\begin{enumerate}
  \item $Th(\Z, S)$ teoria następnika
  \item teoria niezależnych predykatów $Th(2^\omega, P_n)$, $P_n(x)\iff x(n)=0$
  \item teoria niezależnych relacji równoważności o $2$ klasach, $Th(2^\omega, E_n)$, $E_n(x, y)\iff x(n)=y(n)$
  \item teoria z zadania wyżej
  \item teoria nieskończonych przestrzeni liniowych nad ciałem skończonym $F$
  \item teoria przestrzeni liniowych nad przeliczalnym ciałem nieskończonym $F$
\end{enumerate}

\begin{zadanie}
  Dla każdego z przykładów teorii $T$ wyżej opisać zupełne $1$-typy nad modelem $M\models T$ oraz rozstrzygnąć, czy jest ona stabilna, superstabilna, $\aleph_0$-stabilna.
\end{zadanie}

\begin{proof}
  \begin{enumerate}
    \item $S_1(\emptyset)$ ma tylko jeden typ bo średnio cokolwiek można wyrazić jak nie mamy punktu zaczepienia. $S_1(A)$ ma dwa przypadki:
      \begin{itemize}
        \item $p(x)=\{x=S^n(a)\}$ dla jakiegoś $a\in A$ oraz $n\in\Z$ -> takie pyśki powinny być realizowane w każdym modelu tej teorii, czy to liczby całkowite, czy rozszerzone liczby całkowite jak te z analizy niestandardowej
        \item $p(x)=\{x\neq S^n(a)\}$ dla wszystkich $a\in A$ oraz $n\in\Z$ -> ten typ może nie być realizowany, ale jak wezmę skończony podzbiór to jestem w stanie wykombinować $x$ który to spełni
      \end{itemize}

      Czy jest stabilne? Tak! Jeśli $|A|=\kappa$ to typów pierwszego rodzaju mam $|A\times \Z|=\kappa\cdot \aleph_0=\kappa$ jeśli $\kappa\geq \aleph_0$ plus jeden typ "brzydki". Czyli mam $\kappa+1=\kappa$ typów.

      Czy jest $\aleph_0$-stabilne? Tak! Z tego samego powodu co wyżej
      
      Czy jest superstabilne? Tak! Dla $\mu=\aleph_0$.

    \item Tutaj nawet $S_1(\emptyset)$ jest w stanie opisać nam dowolny element $2^\omega$ i dodawanie parametrów nic nie zmienia w sumie. Czyli zawsze mam $2^\aleph_0$ typów. Jeśli dorzucę parametry to tylko dodaję rozróżnienie dla typów zawierających $x=a$ dla $a\in A$, czyli dana maska binarna może pojawić się kilka razy teoretycznie. Czyli $|S_1(A)|\max(|A|+2^{\aleph_0})$.

      Czy jest stabilna? Tak, np. dla $\kappa=2^{\aleph_0}$.

      Czy jest superstabilna? Tak, od $\mu=2^{\aleph_0}$ w górę.
      
      Czy jest $\aleph_0$-stabilna? Za cholerę nie.

    \item $S_1(\emptyset)$ chyba ma tylko jeden typ, bo nie ma na czym zawiesić informacji o elemencie. Co jeśli mam parametry? Wtedy dowolny ciąg $2^\omega$ mogę jednoznacznie wyznaczyć na tle tych parametrów, więc mam znowu $|S_1(A)|=\max(2^\aleph_0,|A|)$, bo dokładam bezkarnie formuły $x=a$ dla $a\in A$. 

      Czyli historia powtarza się...

    \item Znowu $S_1(\emptyset)$ ma jeden typ bo za nic na świecie nie rozróżnię ciągów tak o. Znowu mam dwa rodzaje typów imo:
      \begin{itemize}
        \item typy które na $n$-tej współrzędnej biorą jakiś element z $A$ i przyrównują do niego -> takich mam $|A|^{\aleph_0}$ 
        \item mam wszystkie pozostałe elementy w worku p.t. "$x(n)\neq a(n)$" dla wszystkich $n$ i $a\in A$
      \end{itemize}
      To mi daje coś w stylu $|A|^{\aleph_0}$ typów.
      
      Czy jest $\aleph_0$-stabilna? Nie!

      Czy jest stabilna? Pewnie tak, bo dla dużych $|A|$ będzie $|A|^{\aleph_0}=|A|$

      To też mi mówi, że będzie superstabilne ale za mało umiem teorię mnogości żeby powiedzieć kiedy.

    \item na poprzedniej liście: ta teoria ma redukcję kwantyfikatorów

      Niech $M\models T$. Wtenczas $p(x)\in S_1(M)$ ma formę jednego z poniższych
      \begin{itemize}
        \item $p(x)=\{'x=m'\}$ dla pewnego $m\in M$
        \item $p(x)=\{'x\neq\sum\lambda_im_i'\;:\;\lambda_i\in F_q,\;m_i\in M\}$
      \end{itemize}
      czyli $|S_1(M)|\leq |M|+1=|M|$

      %
      % humpty dumpty $S_1(\emptyset)$ to będzie co, albo zero albo dowolna rzecz? Bo mam równania jednego parametru to mogę je skalować i śmiga. Co jeśli dorzucę parametry? Potrafię zapisać wtedy kombinacje liniowe wszystkich rzeczy z $A$ albo mogę powiedzieć, że mam wektor spoza $span(A)$ bo sobie neguję wszystkie możliwe kombinacje liniowe elementów z $A$ i jestem szczęśliwa. Czyli
      % \begin{itemize}
      %   \item $x=\sum \lambda_i a_i$ dla $\lambda_i\in F$ i $a_i\in A$ -> takich mam $|A\times F|$, co dla nieskończonych $A$ powinno być równe $|A|$
      %   \item ten jeden typ "nie w spanie"
      % \end{itemize}

      Czy jest $\aleph_0$-stabilne? Tak!

      Czy jest stabilne? No jasne

      Czy jest superstabilne? Tak! Wystarczy $\mu=\aleph_0$.
      
    \item tutaj powinno być podobnie. Dalej mam podział na kombinacje liniowe elementów z $A$, czyli mam nie więcej jaik $|A\times F|$ kombinacji liniowych o współczynnikach z $F$

      Czy jest $\aleph_0$-stabilne? Tak

      Czy jest stabilne? Tak!

      Czy jest super stabilne? Tak!
  \end{enumerate}
\end{proof}

\subsection{Modele pierwsze i atomowe}

\begin{definition}{model pierwszy, atomowy}{}
  Model $M\models T$ teorii zupełnej nazywamy
  \begin{description}
    \item[atomowym] gdy $(\forall\;\overline{a}\subseteq M)\;tp(\overline{a})$ jest izolowany (istnieje $\phi$ takie, że $\phi\vdash tp(\overline{a})$)
    \item[pierwszym] gdy $(\forall\;N\models T)\exists\;f:M\to N$ elementarne
  \end{description}
\end{definition}

Model $M$ jest minimalny, gdy nie ma elementarnego podmodelu.

\begin{zadanie}
  Które z teorii wyżej maja model pierwszy? Czy jest on minimalny?
\end{zadanie}

\begin{theorem}{}{}
  Jeśli $T$ jest teorią $\aleph_0$-stabilną to ma model pierwszy.
\end{theorem}

\begin{theorem}{}{}
  Model pierwszy jest jedyny z dokładnością do izomorfizmu.
\end{theorem}

\begin{proof}
  \begin{enumerate}
    \item Ma model pierwszy, bo jest $\aleph_0$-stabilna. Twierdze, że jest nim $(\Z, S)$
      \begin{description}
        \item[pierwszość=atomowość] bierę dowolny $(a_1,...,a_n)\in\Z^n$, jego typ jest izolowany przez $\phi(x_1,...,x_n)="x_1=S^{a_1-a_2}(x_2)\;\land x_1=S^{a_1-a_3}(x_3)\;\land\;..."$
        \item[minimalność] nie wprost: jeśli $N\prec (\Z, S)$ ale $N\neq (\Z, S)$ to mogę wziąć $x\in \Z-N$ oraz $y\in N\subseteq \Z$. Wtedy dla $n=x-y$ zachodzi $x=S^n(y)$, ale to powinno być w $N$ a nie jest.
      \end{description}
    \item Tutaj nie ma modelu pierwszego, bo nie ma typów izolowanych. Każda formuła smyra tylko skończenie wiele indeksów, a żeby wyizolować jakiś typ trzebaby dokładnie określić co się dzieje na wszystkich indeksach.
      
      \textbf{Brak typów izolowanych} załóżmy nie wprost, że $(a_1,...,a_n)\in 2^\omega$ ma typ izolowany przez pewną formułę $\psi(x_1,...,x_n)=\bigwedge P_{n_k}(x_{i_k})\;\land\;\bigwedge \neg P_{n_j}(x_{i_j})$. Niech $N$ będzie największym indeksem występującym w formule $\psi$. Wtedy krotka 
      $$b_i=\begin{cases}a_i & i \leq N\\ -a_i & i > N \end{cases}$$
      również spełnia $\psi$. Czyli $\psi\in tp(a_1,...,a_n)$ oraz $\psi\in tp(b_1,...,b_n)$, co by oznaczało że te typy są takie same, ale nie są - sprzeczność.
    \item skip
    \item skip
    \item Ma model pierwszy, gdyż iż i ponieważ jest $\aleph_0$-stabilne. Zapewne to jest przestrzeń wektorowa o nieskończonej, przeliczalnej bazie - inaczej: ciągi elementów z $F$ o skończenie wielu niezerowych elementach.
      \begin{description}
        \item[pierwszość=atomowość] bierę dowolnego pyśka $(a_1,...,a_n)\in V^n$ -> mogę zapisac formułę $\phi(x_1,...,x_n)$ która będzie konjiunkcją informacji o tym które pyśki są liniowo zależne, a które są liniowo niezależne. Czyli $(\exists \lambda_1...\lambda_i)\;\lambda_i\neq 0 \;\land\;\lambda_1x_{k_1}+...+\lambda_ix_{k_i}=0$ oraz $(\forall\;\lambda_1...\lambda_i)\; \lambda_i\neq 0\;\land\;\lambda_1x_{k_1}+...\lambda_ix_{k_i}\neq 0$
        \item[minimalność] no nie, bo mogę usunąć np. pierwszy wektor bazowy (z nieskończenie wielu), nadal mieć nieskończenie wiele wektorów bazowych i mieć elegancką podprzestrzeń liniową
      \end{description}
    \item \color{red} może nie tak szybko? Ma model pierwszy jeszcze raz, bo jest $\aleph_0$-stabilne. Zapewne znowu to będzie nieskończenie wymiarowa przestrzeń liniowa i argumenty będą bardzo podobne jak wcześniej.
  \end{enumerate}
\end{proof}

\begin{zadanie}
  \begin{enumerate}
    \item Załóżmy, że $T$ ma model pierwszy i nie jest on minimalny. Udowodnić, że wtedy istnieje nieprzeliczalny model atomowy teorii $T$.
    \item Udowodnić także implikację odwrotną.
  \end{enumerate}
\end{zadanie}

\begin{proof}
  \begin{enumerate}
    \item Niech $M$ będzie nieminimalnym modelem pierwszym $T$. Fakt z wykładu:
      \begin{fact}{chyba z konstrukcji ale też z dawida markera}{}
        $M\models T$ jest pierwszy $\iff$ jest przeliczalny i atomowy (część że pierwszy $\implies$ atomowy jest z wykładu).
      \end{fact}
      $M$ nie jest minimalne $\implies$ istnieje $N_0\prec M$, $N_0\neq M$

      $M$ jest pierwsze, czyli istnieje elementarne włożenie $f_0:M\to N_0$

      w modelu monstrum $\mathcal{M}$ widzę to elementarne włożenie jako obcięcie pewnego automorfizmu

      w takim razie $g_0=f_0^{-1}$ jest sensowne i mogę zdefiniować ciąg
      $$M_0=M\prec M_1=g_0(M_0)\prec ...$$
      dla $i<\omega$

      pierwsze przejście graniczne to
      $$M_\omega=\bigcup_{i<\omega} M_i$$
      po tym przejściu granicznym mam model, który jest przeliczalny i typ każdej krotki jest izolowany, bo tak było po drodze -> jest to izomorficzne z typem pierwszym

      czyli teraz mogę znaleźć nowy podmodel $N_1\prec M_\omega$, $N_1\neq M_\omega$ i elementarne włożenie $f_1:M_\omega\to N_1$ które da się odwrócić jako automorfizm modelu monstrum, by skonstruować nowy ciąg -> po $i<\omega_1$, które jest już chyba nieprzeliczalne, a jak nie to mogę powtórzyć i kiedyś osiągnę to nieprzeliczalne piekło.

    \item 
      \begin{description}
        \item[istnienie pierwszego] tak jak w dowodzie istnienia modelu pierwszego dla $\aleph_0$-stabilnej teorii 

          wybieramy $\{a_n\;:\;n<\omega\}$ spełniające test tarskiego-vaughta (czyli będące elementarnym podmodelem) takie, że $tp(a_n/a_{<n})$ jest izolowany (co mogę zapewnić, bo jest to twierdzenie że $tp(a)$ i $tp(b/a)$ są izolowane $\iff$ $tp(a,b)$ jest izolowany)

          to mi daje przeliczalny model atomowy, czyli jest to model pierwszy
        \item[nie jest minimalny] załóżmy nie wprost, że $M_0$ z konstrukcji wyżej jest minimalnym modelem pierwszym


      \end{description}

  \end{enumerate}
\end{proof}



